% !TeX root = resume_ru.tex

\documentclass[10pt, a4paper]{article}

\usepackage[russian]{babel}

\usepackage{trimclip}
\usepackage{fontspec} % Для работы с системными шрифтами (нужен XeLaTeX или LuaLaTeX)
\usepackage{geometry} % Для настройки полей страницы
\usepackage{graphicx} % Для вставки изображений
\usepackage[svgnames]{xcolor} % Для работы с цветами
\usepackage[hidelinks]{hyperref}   % для кликабельных ссылок
\usepackage{enumitem}   % для настройки списков
\usepackage{fontawesome5} % Для иконок (email, телефон и т.д.)
\usepackage{tikz}       % для рисования графических элементов (линия времени, кружки)
\usetikzlibrary{shapes.misc, positioning, arrows.meta}
\usepackage{changepage}

% --- НАСТРОЙКИ ШАБЛОНА ---

% 1. ГЕОМЕТРИЯ СТРАНИЦЫ
\geometry{
    a4paper,
    left=12mm,
    right=12mm,
    top=10mm,
    bottom=10mm,
    footskip=.5cm
}

% 2. ЦВЕТА
\definecolor{maintext}{HTML}{333333} % Основной цвет текста
\definecolor{graytext}{HTML}{666666} % Серый цвет для дат и подписей
\definecolor{boxgray}{HTML}{EBEBEB}  % Цвет для "таблеток" в интересах

% 3. ШРИФТЫ
\setmainfont{Lato} % Используем шрифт Lato
\color{maintext}
\pagestyle{empty} % Убираем нумерацию страниц

% 4. НАСТРОЙКИ ССЫЛОК
\hypersetup{
    colorlinks=true,
    urlcolor=maintext,
    linkcolor=maintext,
    pdfborder={0 0 0}, % убираем рамку вокруг ссылок
}

\setlength{\parindent}{0pt}

\setlist[itemize]{nosep, topsep=2pt, itemsep=1.5pt, leftmargin=*}
\linespread{0.94}

\begin{document}
% !TeX root = ../resume.tex

% --- ВСПОМОГАТЕЛЬНЫЕ КОМАНДЫ ---

% Команда для заголовков секций (WORK EXPERIENCE, SKILLS и т.д.)
\newcommand{\sectiontitle}[1]{%
    \vspace{8pt}
    \par
    \noindent\MakeUppercase{#1}
    \hrule width 1\textwidth height 0.4pt
    \vspace{8pt}
}

% Команда для иконки внешней ссылки
\newcommand{\extlink}{\, \textcolor{graytext}{\faIcon{external-link-alt}}}

% Команда для пунктов в списке достижений
\setlist[itemize,1]{label=\,--, leftmargin=*, topsep=0pt, itemsep=4pt}

% Команда для описания языков
\newcommand{\languageskill}[2]{%
    \par\noindent\hangindent=1em\hangafter=0
    \makebox[8em][l]{#1} % Название языка
    \foreach \i in {1,...,5}{
        \ifnum\i > #2
            \tikz\draw[graytext, fill=white] (0,0) circle (2.5pt);
        \else
            \tikz\fill[maintext] (0,0) circle (2.5pt);
        \fi
        \hspace{2pt}
    }
    \vspace{4pt}
}

% Команда для "таблеток" с интересами
\newcommand{\interesttag}[1]{%
    \tikz[baseline=(X.base)] \node[draw=graytext, rounded corners=3pt, inner sep=4pt, text=graytext] (X) {#1};
}

% ===============================================================
% ШАПКА (ИМЯ, КОНТАКТЫ)
% ===============================================================

\begin{minipage}[t]{0.6\textwidth} % Левая часть шапки
    \hspace{0.02\textwidth} % Отступ между краем и текстом
    \begin{minipage}[c]{0.7\textwidth}
        \Huge \textbf{Сергей Новиков}
    \end{minipage}
\end{minipage}%
\begin{minipage}{0.4\textwidth} % Правая часть шапки
    \raggedleft
    \small
    \href{tel:+79635716695}{+7 (963) 571-66-95} \faIcon{mobile-alt} \\[4pt]
    \href{mailto:isnov.contact@gmail.com}{isnov.contact@gmail.com} \faEnvelope \\[4pt]
    \href{https://t.me/Isn0v}{@Isn0v} \faIcon{telegram-plane} \\[4pt]
    \href{https://github.com/Isn0v}{github.com/Isn0v} \faIcon{github} \\[4pt]
    \href{https://leetcode.com/u/Isn0v/}{leetcode.com/Isn0v} \faIcon{code} \\[4pt]
\end{minipage}

\vspace{0.8cm}

% ===============================================================
% ДВУХКОЛОНОЧНАЯ СТРУКТУРА
% ===============================================================

\begin{minipage}[t]{0.55\textwidth}

\raggedright

% --- ОПЫТ РАБОТЫ ---
\sectiontitle{\faBriefcase \quad Опыт работы}

\begin{tikzpicture}[remember picture]
  \node[circle, fill=maintext, inner sep=2.5pt] (c1) at (0,0) {};
\end{tikzpicture}
\hspace{0.5cm}
\begin{minipage}[t]{0.9\textwidth}
  \raggedright
  \textbf{\Large SRE-инженер} \\[2pt]
  \textcolor{graytext}{\href{https://www.sberbank.ru/ru/person}{СберБанк}}\extlink \\
  \textcolor{graytext}{\textit{с 11/2025}}
  \begin{itemize}
      \item \textbf{Разработал и интегрировал} со стеком наблюдаемости (ELK/Prometheus) AI-агента для анализа логов, ускорив поиск первопричин (RCA) и сократив время триажа инцидентов на ~40\%.
      \item \textbf{Спроектировал и стандартизировал} Helm-чарты и CI/CD-пайплайны деплоя критически важных enterprise-сервисов (GitLab, Jira, Confluence, Nexus) в K8s/OpenShift (доступность 99.9\% для 3\,000+ пользователей).
      \item \textbf{Провел рефакторинг} кодовой базы Helm-чартов Nextcloud: сократил объем конфигурации на 60\% (с 25k до 10k строк) с сохранением полной функциональности и устранением дрейфа конфигураций.
      \item \textbf{Оптимизировал} конфигурацию NGINX Ingress / балансировщиков и разработал патчи для Harbor, устранив узкие места и повысив стабильность выкачки артефактов под нагрузкой.
  \end{itemize}
\end{minipage}

\vspace{0.7cm}

\begin{tikzpicture}[remember picture]
    \node[circle, fill=maintext, inner sep=2.5pt] (c1) at (0,0) {};
\end{tikzpicture}
\hspace{0.5cm}
\begin{minipage}[t]{0.9\textwidth}
    \raggedright
    \textbf{\Large DevOps-инженер} \\[2pt]
    \textcolor{graytext}{\href{https://rnew.tilda.ws/excelsiorathuawei}{Excelsior at Huawei}}\extlink \\
    \textcolor{graytext}{\textit{03/2025 -- 07/2025}}
    \begin{itemize}
        \item \textbf{Модульно переработал} GitLab CI/CD пайплайны системной разработки, сократив среднее время сборки и тестов на ~30\% за счет параллелизации и кэширования артефактов.
        \item \textbf{Внедрил} шаблоны пайплайнов, ускорив доставку изменений для ключевой команды разработки.
    \end{itemize}
\end{minipage}

\vspace{0.7cm}

\begin{tikzpicture}[remember picture]
    \node[circle, fill=maintext, inner sep=2.5pt] (c1) at (0,0) {};
\end{tikzpicture}
\hspace{0.5cm}
\begin{minipage}[t]{0.9\textwidth}
    \raggedright
    \textbf{\Large Стажер-исследователь Linux} \\[2pt]
    \textcolor{graytext}{\href{https://rnew.tilda.ws/excelsiorathuawei}{Excelsior at Huawei}}\extlink \\
    \textcolor{graytext}{\textit{09/2024 -- 01/2025}}
    \begin{itemize}
        \item \textbf{Исследовал} внутренние механизмы ядра Linux: планировщик задач (CFS), виртуальную файловую систему (VFS), системные вызовы и управление памятью процессов.
        \item \textbf{Проводил бенчмаркинг} примитивов изоляции ресурсов и контейнеризации ядра (namespaces, cgroups, seccomp) для оценки накладных расходов и производительности.
    \end{itemize}
\end{minipage}

\end{minipage}%
\hfill % Разделитель между колонками
\begin{minipage}[t]{0.4\textwidth}

\sectiontitle{\faFolderOpen \quad Проекты}

\begin{tikzpicture}[remember picture]
    \node[circle, fill=maintext, inner sep=2.5pt] (c1) at (0,0) {};
\end{tikzpicture}
\hspace{0.5cm}
\begin{minipage}[t]{0.9\textwidth}
    \raggedright
    \textbf{\Large \href{https://github.com/Isn0v/Sentix}{Sentix}\extlink}  \\[2pt]
    Микросервисный пайплайн на Go (конкурентный сбор) и Python для анализа сентимента крипто- и финансовых новостей в реальном времени (OpenAI API, PostgreSQL, Docker Compose). \\
\end{minipage}

\vspace{0.4cm}

\begin{tikzpicture}[remember picture]
    \node[circle, fill=maintext, inner sep=2.5pt] (c1) at (0,0) {};
\end{tikzpicture}
\hspace{0.5cm}
\begin{minipage}[t]{0.9\textwidth}
    \raggedright
    \textbf{\Large \href{https://github.com/B0GDANPN/BrickOS/tree/main}{BrickOS}\extlink}  \\[2pt]
    Собственное ядро x86-совместимой ОС с нуля на C и x86 Assembly: обработка аппаратных прерываний и управление памятью.
\end{minipage}


\sectiontitle{\faGraduationCap \quad Образование}
\begin{minipage}[t]{\textwidth}
    \raggedright
    \textbf{\Large Математика и \\компьютерные науки} \\[2pt]
    \textit{\large \href{https://education.nsu.ru/syspro/}{Системное программирование} \extlink} \\[4pt]
    \textcolor{graytext}{Новосибирский госуниверситет (НГУ)} \\
    \textcolor{graytext}{\textit{Механико-математический факультет, \\бакалавриат}\\ \textit{2022 -- 2026}}
\end{minipage}

\vspace{0.4cm}

% --- Исследования ---

\sectiontitle{\faBookOpen \quad Научная деятельность}
\begin{minipage}[t]{\textwidth}
  \raggedright
  \large \href{https://www.nsu.ru/upload/iblock/826/371ay2z05ae8pdhbgjr5b14hcokg31t1/%D0%98%D0%A2_%D0%9D%D0%B0%D1%83%D1%87%D0%BD%D1%8B%D0%B9%20%D0%B8%D0%BD%D0%B6%D0%B8%D0%BD%D0%B8%D1%80%D0%B8%D0%BD%D0%B3.pdf}{\textbf{Моделирование когнитивных агентов нейроэволюцией}}\extlink  \\[4pt]
  \small \textcolor{graytext}{Публикация на МНСК-2026 (с. 168). Продемонстрировал преимущество NEAT перед DQN в зашумленной среде.} \\
\end{minipage}

\vspace{0.4cm}

% --- НАВЫКИ ---
\sectiontitle{\faCog \quad Навыки}
\begin{minipage}[t]{1\textwidth}
    \raggedright
    \small
    Linux \\[4pt]
    Алгоритмы и структуры данных \\[4pt]
    GitLab CI/CD \\[4pt]
    Git \& GitHub \\[4pt]
    C, C++, Go, Python \\[4pt]
    Компьютерные сети \\[4pt]
    Docker, K8s, OpenShift, Helm, Kustomize, NGINX \\[4pt]
    PostgreSQL \\[4pt]
    Grafana, Prometheus, ELK
\end{minipage}

\vspace{0.4cm}

% --- ЯЗЫКИ ---
\sectiontitle{\faGlobe \quad Языки}
Русский --- родной \\[4pt]
Английский --- B2


\end{minipage}

\end{document}
